\section{Bases and sheaves}
\label{section-bases}

\noindent
Sometimes there exists a basis for the topology
consisting of opens that are easier to work with
than general opens. For convenience we give here
some definitions and simple lemmas in order to
facilitate working with (pre)sheaves in such a situation.

\begin{definition}
\label{definition-presheaf-basis}
Let $X$ be a topological space. Let $\mathcal{B}$ be a
basis for the topology on $X$.
\begin{enumerate}
\item A {\it presheaf $\mathcal{F}$ of sets on $\mathcal{B}$}
is a rule which assigns to each $U \in \mathcal{B}$ a set
$\mathcal{F}(U)$ and to each inclusion $V \subset U$
of elements of $\mathcal{B}$ a map
$\rho^U_V : \mathcal{F}(U) \to \mathcal{F}(V)$ such that
$\rho^U_U = \text{id}_{\mathcal{F}(U)}$ for all $U \in \mathcal{B}$
whenever $W \subset V \subset U$ in $\mathcal{B}$ we have
$\rho^U_W = \rho^V_W \circ \rho ^U_V$.
\item A {\it morphism $\varphi : \mathcal{F} \to \mathcal{G}$
of presheaves of sets on $\mathcal{B}$} is a rule which assigns to each
element $U \in \mathcal{B}$ a map of sets $\varphi : \mathcal{F}(U)
\to \mathcal{G}(U)$ compatible with restriction maps.
\end{enumerate}
\end{definition}

\noindent
As in the case of usual presheaves we use the terminology of sections,
restrictions of sections, etc. In particular, we may define the
{\it stalk} of $\mathcal{F}$ at a point $x \in X$ by the
colimit
$$
\mathcal{F}_x = \colim_{U\in \mathcal{B}, x\in U} \mathcal{F}(U).
$$
As in the case of the stalk of a presheaf on $X$ this limit is
directed. The reason is that the collection of $U\in \mathcal{B}$,
$x \in U$ is a fundamental system of open neighbourhoods of $x$.

\medskip\noindent
It is easy to make examples to show that the notion of a presheaf
on $X$ is very different from the notion of a presheaf on a basis
for the topology on $X$. This does not happen in the case of
sheaves. A much more useful notion therefore, is the following.

\begin{definition}
\label{definition-sheaf-basis}
Let $X$ be a topological space. Let $\mathcal{B}$ be a
basis for the topology on $X$.
\begin{enumerate}
\item A {\it sheaf $\mathcal{F}$ of sets on $\mathcal{B}$} is a presheaf
of sets on $\mathcal{B}$ which satisfies the following additional
property: Given any $U \in \mathcal{B}$, and any covering
$U = \bigcup_{i \in I} U_i$ with $U_i \in \mathcal{B}$, and
any coverings $U_i \cap U_j = \bigcup_{k \in I_{ij}} U_{ijk}$ with
$U_{ijk} \in \mathcal{B}$ the sheaf condition holds:
\begin{itemize}
\item[(**)] For any collection
of sections $s_i \in \mathcal{F}(U_i)$, $i \in I$ such that
$\forall i, j\in I$, $\forall k\in I_{ij}$
$$
s_i|_{U_{ijk}} = s_j|_{U_{ijk}}
$$
there exists a unique section $s \in \mathcal{F}(U)$ such that
$s_i = s|_{U_i}$ for all $i \in I$.
\end{itemize}
\item A {\it morphism of sheaves of sets on $\mathcal{B}$} is simply a
morphism of presheaves of sets.
\end{enumerate}
\end{definition}

\noindent
First we explain that it suffices to
check the sheaf condition $(**)$
on a cofinal system of coverings.
In the situation of the definition, suppose
$U \in \mathcal{B}$. Let us temporarily denote
$\text{Cov}_\mathcal{B}(U)$ the set of all coverings
of $U$ by elements of $\mathcal{B}$.
Note that $\text{Cov}_\mathcal{B}(U)$ is preordered by refinement.
A subset $C \subset \text{Cov}_\mathcal{B}(U)$ is a
cofinal system, if for every $\mathcal{U} \in \text{Cov}_\mathcal{B}(U)$
there exists a covering $\mathcal{V} \in C$ which refines $\mathcal{U}$.

\begin{lemma}
\label{lemma-cofinal-systems-coverings}
With notation as above.
For each $U \in \mathcal{B}$, let $C(U) \subset \text{Cov}_\mathcal{B}(U)$
be a cofinal system. For each $U \in \mathcal{B}$, and each
$\mathcal{U} : U = \bigcup U_i$ in $C(U)$, let coverings
$\mathcal{U}_{ij} : U_i \cap U_j = \bigcup U_{ijk}$,
$U_{ijk} \in \mathcal{B}$ be given.
Let $\mathcal{F}$ be a presheaf of sets on $\mathcal{B}$.
The following are equivalent
\begin{enumerate}
\item The presheaf $\mathcal{F}$ is a sheaf on $\mathcal{B}$.
\item For every $U \in \mathcal{B}$ and every covering
$\mathcal{U} : U = \bigcup U_i$ in $C(U)$ the sheaf condition
$(**)$ holds (for the given coverings $\mathcal{U}_{ij}$).
\end{enumerate}
\end{lemma}

\begin{proof}
We have to show that (2) implies (1).
Suppose that $U \in \mathcal{B}$, and that
$\mathcal{U} : U = \bigcup_{i\in I} U_i$ is an arbitrary covering
by elements of $\mathcal{B}$.  Because the system $C(U)$ is cofinal
we can find an element $\mathcal{V} : U = \bigcup_{j \in J} V_j$
in $C(U)$ which refines $\mathcal{U}$. This means there exists
a map $\alpha : J \to I$ such that $V_j \subset U_{\alpha(j)}$.

\medskip\noindent
Note that if $s, s' \in \mathcal{F}(U)$ are sections such
that $s|_{U_i} = s'|_{U_i}$, then
$$
s|_{V_j}
= (s|_{U_{\alpha(j)}})|_{V_j}
= (s'|_{U_{\alpha(j)}})|_{V_j}
= s'|_{V_j}
$$
for all $j$. Hence by the uniqueness in $(**)$
for the covering $\mathcal{V}$ we conclude that $s = s'$.
Thus we have proved the uniqueness part of $(**)$
for our arbitrary covering $\mathcal{U}$.

\medskip\noindent
Suppose furthermore that $U_i \cap U_{i'} = \bigcup_{k \in I_{ii'}} U_{ii'k}$
are arbitrary coverings by $U_{ii'k} \in \mathcal{B}$.
Let us try to prove the existence part of $(**)$ for the system
$(\mathcal{U}, \mathcal{U}_{ij})$. Thus let $s_i \in \mathcal{F}(U_i)$
and suppose we have
$$
s_i|_{U_{ii'k}} = s_{i'}|_{U_{ii'k}}
$$
for all $i, i', k$. Set $t_j = s_{\alpha(j)}|_{V_j}$, where $\mathcal{V}$
and $\alpha$ are as above.

\medskip\noindent
There is one small kink in the argument here. Namely, let
$\mathcal{V}_{jj'} : V_j \cap V_{j'} = \bigcup_{l \in J_{jj'}} V_{jj'l}$
be the covering given to us by the statement of the lemma.
It is not a priori clear that
$$
t_j|_{V_{jj'l}} = t_{j'}|_{V_{jj'l}}
$$
for all $j, j', l$. To see this, note that we do have
$$
t_j|_W = t_{j'}|_W \text{ for all } W \in \mathcal{B},
W \subset V_{jj'l} \cap U_{\alpha(j)\alpha(j')k}
$$
for all $k \in I_{\alpha(j)\alpha(j')}$, by our assumption on
the family of elements $s_i$. And since
$V_j \cap V_{j'} \subset U_{\alpha(j)} \cap U_{\alpha(j')}$
we see that $t_j|_{V_{jj'l}}$ and $t_{j'}|_{V_{jj'l}}$
agree on the members of a covering of $V_{jj'l}$ by
elements of $\mathcal{B}$. Hence by the uniqueness part proved above
we finally deduce the desired equality of
$t_j|_{V_{jj'l}}$ and $t_{j'}|_{V_{jj'l}}$.
Then we get the existence of an element $t \in \mathcal{F}(U)$
by property $(**)$ for $(\mathcal{V}, \mathcal{V}_{jj'})$.

\medskip\noindent
Again there is a small snag. We know that $t$ restricts to $t_j$ on $V_j$
but we do not yet know that $t$ restricts to $s_i$ on $U_i$. To conclude
this note that the sets $U_i \cap V_j$, $j \in J$ cover $U_i$. Hence
also the sets $U_{i \alpha(j) k} \cap V_j$, $j\in J$, $k \in I_{i\alpha(j)}$
cover $U_i$. We leave it to the reader to see that $t$ and $s_i$ restrict
to the same section of $\mathcal{F}$ on any $W \in \mathcal{B}$
which is contained in one of the open sets
$U_{i \alpha(j) k} \cap V_j$, $j\in J$, $k \in I_{i\alpha(j)}$.
Hence by the uniqueness part seen above we win.
\end{proof}

\begin{lemma}
\label{lemma-cofinal-systems-coverings-standard-case}
Let $X$ be a topological space. Let $\mathcal{B}$ be a basis for the
topology on $X$. Assume that for every triple $U, U', U'' \in \mathcal{B}$
with $U' \subset U$ and $U'' \subset U$ we have $U' \cap U'' \in \mathcal{B}$.
For each $U \in \mathcal{B}$, let $C(U) \subset \text{Cov}_\mathcal{B}(U)$
be a cofinal system.
Let $\mathcal{F}$ be a presheaf of sets on $\mathcal{B}$.
The following are equivalent
\begin{enumerate}
\item The presheaf $\mathcal{F}$ is a sheaf on $\mathcal{B}$.
\item For every $U \in \mathcal{B}$ and every covering
$\mathcal{U} : U = \bigcup U_i$ in $C(U)$ and for every
family of sections $s_i \in \mathcal{F}(U_i)$ such
that $s_i|_{U_i \cap U_j} = s_j|_{U_i \cap U_j}$ there
exists a unique section $s \in \mathcal{F}(U)$ which
restricts to $s_i$ on $U_i$.
\end{enumerate}
\end{lemma}

\begin{proof}
This is a reformulation of
Lemma \ref{lemma-cofinal-systems-coverings} above
in the special case where the coverings $\mathcal{U}_{ij}$
each consist of a single element. But also this case is much
easier and is an easy exercise to do directly.
\end{proof}

\begin{lemma}
\label{lemma-condition-star-sections}
Let $X$ be a topological space.
Let $\mathcal{B}$ be a basis for the topology on $X$.
Let $U \in \mathcal{B}$.
Let $\mathcal{F}$ be a sheaf of sets on $\mathcal{B}$.
The map
$$
\mathcal{F}(U) \to \prod\nolimits_{x \in U} \mathcal{F}_x
$$
identifies $\mathcal{F}(U)$ with the elements $(s_x)_{x\in U}$
with the property
\begin{itemize}
\item[(*)] For any $x \in U$ there exists a $V \in \mathcal{B}$,
with $x \in V \subset U$ and a section $\sigma \in \mathcal{F}(V)$
such that for all $y \in V$ we have $s_y = (V, \sigma)$ in $\mathcal{F}_y$.
\end{itemize}
\end{lemma}

\begin{proof}
First note that the map
$\mathcal{F}(U) \to \prod\nolimits_{x \in U} \mathcal{F}_x$
is injective by the uniqueness in the sheaf condition
of Definition \ref{definition-sheaf-basis}. Let $(s_x)$ be
any element on the right hand side which satisfies $(*)$.
Clearly this means we can find a covering $U = \bigcup U_i$,
$U_i \in \mathcal{B}$ such that $(s_x)_{x \in U_i}$ comes from
certain $\sigma_i \in \mathcal{F}(U_i)$. For every $y \in U_i \cap U_j$
the sections $\sigma_i$ and $\sigma_j$ agree in the stalk
$\mathcal{F}_y$. Hence there exists an element $V_{ijy} \in \mathcal{B}$,
$y \in V_{ijy}$ such that $\sigma_i|_{V_{ijy}} = \sigma_j|_{V_{ijy}}$.
Thus the sheaf condition $(**)$ of Definition \ref{definition-sheaf-basis}
applies to the system of $\sigma_i$ and we obtain a section
$s \in \mathcal{F}(U)$ with the desired property.
\end{proof}

\noindent
Let $X$ be a topological space.
Let $\mathcal{B}$ be a basis for the topology on $X$.
There is a natural restriction functor from the category
of sheaves of sets on $X$ to the category of sheaves of
sets on $\mathcal{B}$. It turns out that this is an equivalence
of categories. In down to earth terms this means the following.

\begin{lemma}
\label{lemma-extend-off-basis}
Let $X$ be a topological space.
Let $\mathcal{B}$ be a basis for the topology on $X$.
Let $\mathcal{F}$ be a sheaf of sets on $\mathcal{B}$.
There exists a unique sheaf of sets $\mathcal{F}^{ext}$
on $X$ such that $\mathcal{F}^{ext}(U) = \mathcal{F}(U)$
for all $U \in \mathcal{B}$ compatibly with the restriction
mappings.
\end{lemma}

\begin{proof}
We first construct a presheaf $\mathcal{F}^{ext}$ with the
desired property. Namely, for an arbitrary open $U \subset X$ we
define $\mathcal{F}^{ext}(U)$ as the set of elements
$(s_x)_{x \in U}$ such that $(*)$ of
Lemma \ref{lemma-condition-star-sections} holds.
It is clear that there are restriction mappings
that turn $\mathcal{F}^{ext}$ into a presheaf of sets.
Also, by Lemma \ref{lemma-condition-star-sections} we
see that $\mathcal{F}(U) = \mathcal{F}^{ext}(U)$
whenever $U$ is an element of the basis $\mathcal{B}$.
To see $\mathcal{F}^{ext}$ is a sheaf one may
argue as in the proof of Lemma \ref{lemma-sheafification-sheaf}.
\end{proof}

\noindent
Note that we have
$$
\mathcal{F}_x = \mathcal{F}_x^{ext}
$$
in the situation of the lemma. This is so because the
collection of elements of $\mathcal{B}$ containing
$x$ forms a fundamental system of open neighbourhoods of $x$.

\begin{lemma}
\label{lemma-restrict-basis-equivalence}
Let $X$ be a topological space.
Let $\mathcal{B}$ be a basis for the topology on $X$.
Denote $\Sh(\mathcal{B})$ the category of
sheaves on $\mathcal{B}$.
There is an equivalence of categories
$$
\Sh(X) \longrightarrow \Sh(\mathcal{B})
$$
which assigns to a sheaf on $X$ its restriction to
the members of $\mathcal{B}$.
\end{lemma}

\begin{proof}
The inverse functor in given in Lemma \ref{lemma-extend-off-basis} above.
Checking the obvious functorialities is left to the
reader.
\end{proof}

\noindent
This ends the discussion of sheaves of sets on
a basis $\mathcal{B}$. Let $(\mathcal{C}, F)$ be
a type of algebraic structure. At the end of this section
we would like to point out that the constructions
above work for sheaves with values in $\mathcal{C}$.
Let us briefly define the relevant notions.

\begin{definition}
\label{definition-sheaf-structures-basis}
Let $X$ be a topological space. Let $\mathcal{B}$ be a
basis for the topology on $X$. Let $(\mathcal{C}, F)$ be
a type of algebraic structure.
\begin{enumerate}
\item A {\it presheaf $\mathcal{F}$ with values in $\mathcal{C}$
on $\mathcal{B}$} is a rule which assigns to each
$U \in \mathcal{B}$ an object
$\mathcal{F}(U)$ of $\mathcal{C}$ and to each inclusion $V \subset U$
of elements of $\mathcal{B}$ a morphism
$\rho^U_V : \mathcal{F}(U) \to \mathcal{F}(V)$ in $\mathcal{C}$ such that
$\rho^U_U = \text{id}_{\mathcal{F}(U)}$ for all $U \in \mathcal{B}$ and
whenever $W \subset V \subset U$ in $\mathcal{B}$ we have
$\rho^U_W = \rho^V_W \circ \rho ^U_V$.
\item A {\it morphism $\varphi : \mathcal{F} \to \mathcal{G}$
of presheaves with values in $\mathcal{C}$
on $\mathcal{B}$} is a rule which assigns to each
element $U \in \mathcal{B}$ a morphism of
algebraic structures $\varphi : \mathcal{F}(U) \to \mathcal{G}(U)$
compatible with restriction maps.
\item Given a presheaf $\mathcal{F}$ with values in $\mathcal{C}$
on $\mathcal{B}$ we say that $U \mapsto F(\mathcal{F}(U))$ is the
underlying presheaf of sets.
\item A {\it sheaf $\mathcal{F}$ with values in $\mathcal{C}$
on $\mathcal{B}$} is a presheaf with values in $\mathcal{C}$
on $\mathcal{B}$ whose underlying presheaf of sets is a sheaf.
\end{enumerate}
\end{definition}

\noindent
At this point we can define the {\it stalk} at $x \in X$
of a presheaf with values in $\mathcal{C}$ on $\mathcal{B}$
as the directed colimit
$$
\mathcal{F}_x = \colim_{U\in \mathcal{B}, x\in U} \mathcal{F}(U).
$$
It exists as an object of $\mathcal{C}$
because of our assumptions on $\mathcal{C}$.
Also, we see that the underlying set of $\mathcal{F}_x$
is the stalk of the underlying presheaf of sets on $\mathcal{B}$.

\medskip\noindent
Note that Lemmas \ref{lemma-cofinal-systems-coverings},
\ref{lemma-cofinal-systems-coverings-standard-case} and
\ref{lemma-condition-star-sections} refer to the sheaf
property which we have defined in terms of the associated presheaf
of sets. Hence they generalize without change to the notion
of a presheaf with values in $\mathcal{C}$. The analogue of
Lemma \ref{lemma-extend-off-basis} need some care. Here it is.

\begin{lemma}
\label{lemma-extend-off-basis-structures}
Let $X$ be a topological space. Let $(\mathcal{C}, F)$ be
a type of algebraic structure.
Let $\mathcal{B}$ be a basis for the topology on $X$.
Let $\mathcal{F}$ be a sheaf with values in $\mathcal{C}$
on $\mathcal{B}$.
There exists a unique sheaf $\mathcal{F}^{ext}$ with values in $\mathcal{C}$
on $X$ such that $\mathcal{F}^{ext}(U) = \mathcal{F}(U)$
for all $U \in \mathcal{B}$ compatibly with the restriction
mappings.
\end{lemma}

\begin{proof}
By the conditions imposed on the pair $(\mathcal{C}, F)$ it
suffices to come up with a presheaf $\mathcal{F}^{ext}$
which does the correct thing on the level of underlying
presheaves of sets. Thus our first task is to construct
a suitable object $\mathcal{F}^{ext}(U)$ for all open $U \subset X$.
We could do this by imitating
Lemma \ref{lemma-diagram-fibre-product} in the setting
of presheaves on $\mathcal{B}$. However, a slightly different method
(but basically equivalent) is the following:
Define it as the directed colimit
$$
\mathcal{F}^{ext}(U)
:=
\colim_\mathcal{U} FIB(\mathcal{U})
$$
over all coverings
$\mathcal{U} : U = \bigcup_{i\in I} U_i$ by $U_i \in \mathcal{B}$
of the fibre product
$$
\xymatrix{
FIB(\mathcal{U}) \ar[r] \ar[d] &
\prod\nolimits_{x\in U} \mathcal{F}_x \ar[d] \\
\prod\nolimits_{i\in I} \mathcal{F}(U_i) \ar[r] &
\prod\nolimits_{i \in I} \prod\nolimits_{x\in U_i} \mathcal{F}_x
}
$$
By the usual arguments, see Lemma \ref{lemma-image-contained-in}
and Example \ref{example-application-lemma-image-contained-in}
it suffices to show that this construction on underlying
sets is the same as the definition using $(**)$ above.
Details left to the reader.
\end{proof}

\noindent
Note that we have
$$
\mathcal{F}_x = \mathcal{F}_x^{ext}
$$
as objects in $\mathcal{C}$
in the situation of the lemma. This is so because the
collection of elements of $\mathcal{B}$ containing
$x$ forms a fundamental system of open neighbourhoods of $x$.

\begin{lemma}
\label{lemma-restrict-basis-equivalence-structures}
Let $X$ be a topological space.
Let $\mathcal{B}$ be a basis for the topology on $X$.
Let $(\mathcal{C}, F)$ be a type of algebraic structure.
Denote $\Sh(\mathcal{B}, \mathcal{C})$ the category of
sheaves with values in $\mathcal{C}$ on $\mathcal{B}$.
There is an equivalence of categories
$$
\Sh(X, \mathcal{C})
\longrightarrow
\Sh(\mathcal{B}, \mathcal{C})
$$
which assigns to a sheaf on $X$ its restriction to
the members of $\mathcal{B}$.
\end{lemma}

\begin{proof}
The inverse functor in given in
Lemma \ref{lemma-extend-off-basis-structures} above.
Checking the obvious functorialities is left to the
reader.
\end{proof}

\noindent
Finally we come to the case of (pre)sheaves of modules
on a basis. We will use the easy fact that the category
of presheaves of sets on a basis has products and that
they are described by taking products of values on
elements of the bases.

\begin{definition}
\label{definition-sheaf-modules-basis}
Let $X$ be a topological space. Let $\mathcal{B}$ be a
basis for the topology on $X$. Let $\mathcal{O}$ be
a presheaf of rings on $\mathcal{B}$.
\begin{enumerate}
\item A {\it presheaf of $\mathcal{O}$-modules $\mathcal{F}$
on $\mathcal{B}$} is a presheaf of abelian groups on
$\mathcal{B}$ together with a morphism of presheaves
of sets $\mathcal{O} \times \mathcal{F} \to \mathcal{F}$
such that for all $U \in \mathcal{B}$ the map
$\mathcal{O}(U) \times \mathcal{F}(U) \to \mathcal{F}(U)$
turns the group $\mathcal{F}(U)$ into an $\mathcal{O}(U)$-module.
\item A {\it morphism $\varphi : \mathcal{F} \to \mathcal{G}$
of presheaves of $\mathcal{O}$-modules on $\mathcal{B}$}
is a morphism of abelian presheaves on $\mathcal{B}$
which induces an $\mathcal{O}(U)$-module homomorphism
$\mathcal{F}(U) \to \mathcal{G}(U)$ for every $U \in \mathcal{B}$.
\item Suppose that $\mathcal{O}$ is a sheaf of rings
on $\mathcal{B}$. A {\it sheaf $\mathcal{F}$ of $\mathcal{O}$-modules
on $\mathcal{B}$} is a presheaf of $\mathcal{O}$-modules
on $\mathcal{B}$ whose underlying presheaf of abelian groups
is a sheaf.
\end{enumerate}
\end{definition}

\noindent
We can define the {\it stalk} at $x \in X$
of a presheaf of $\mathcal{O}$-modules on $\mathcal{B}$
as the directed colimit
$$
\mathcal{F}_x = \colim_{U\in \mathcal{B}, x\in U} \mathcal{F}(U).
$$
It is a $\mathcal{O}_x$-module.

\medskip\noindent
Note that Lemmas \ref{lemma-cofinal-systems-coverings},
\ref{lemma-cofinal-systems-coverings-standard-case} and
\ref{lemma-condition-star-sections} refer to the sheaf
property which we have defined in terms of the associated presheaf
of sets. Hence they generalize without change to the notion
of a presheaf of $\mathcal{O}$-modules. The analogue of
Lemma \ref{lemma-extend-off-basis} is as follows.

\begin{lemma}
\label{lemma-extend-off-basis-module}
Let $X$ be a topological space.
Let $\mathcal{B}$ be a basis for the topology on $X$.
Let $\mathcal{O}$ be a sheaf of rings on $\mathcal{B}$.
Let $\mathcal{F}$ be a sheaf of $\mathcal{O}$-modules
on $\mathcal{B}$. Let $\mathcal{O}^{ext}$ be the sheaf
of rings on $X$ extending $\mathcal{O}$ and let
$\mathcal{F}^{ext}$ be the abelian sheaf on $X$ extending
$\mathcal{F}$, see Lemma \ref{lemma-extend-off-basis-structures}.
There exists a canonical map
$$
\mathcal{O}^{ext} \times \mathcal{F}^{ext}
\longrightarrow
\mathcal{F}^{ext}
$$
which agrees with the given map over elements of $\mathcal{B}$
and which endows $\mathcal{F}^{ext}$ with the structure
of an $\mathcal{O}^{ext}$-module.
\end{lemma}

\begin{proof}
It suffices to construct the multiplication map
on the level of presheaves of sets. Perhaps the easiest
way to see this is to prove directly that if
$(f_x)_{x \in U}$, $f_x \in \mathcal{O}_x$
and
$(m_x)_{x \in U}$, $m_x \in \mathcal{F}_x$
satisfy $(*)$, then the element
$(f_xm_x)_{x \in U}$ also satisfies $(*)$.
Then we get the desired result, because in the proof
of Lemma \ref{lemma-extend-off-basis} we construct the extension
in terms of families of elements of stalks satisfying $(*)$.
\end{proof}

\noindent
Note that we have
$$
\mathcal{F}_x = \mathcal{F}_x^{ext}
$$
as $\mathcal{O}_x$-modules in the situation of the lemma.
This is so because the collection of elements of $\mathcal{B}$ containing
$x$ forms a fundamental system of open neighbourhoods of $x$, or
simply because it is true on the underlying sets.

\begin{lemma}
\label{lemma-restrict-basis-equivalence-modules}
Let $X$ be a topological space.
Let $\mathcal{B}$ be a basis for the topology on $X$.
Let $\mathcal{O}$ be a sheaf of rings on $X$.
Denote $\textit{Mod}(\mathcal{O}|_\mathcal{B})$ the category of
sheaves of $\mathcal{O}|_\mathcal{B}$-modules on $\mathcal{B}$.
There is an equivalence of categories
$$
\textit{Mod}(\mathcal{O})
\longrightarrow
\textit{Mod}(\mathcal{O}|_\mathcal{B})
$$
which assigns to a sheaf of $\mathcal{O}$-modules on $X$ its restriction to
the members of $\mathcal{B}$.
\end{lemma}

\begin{proof}
The inverse functor in given in
Lemma \ref{lemma-extend-off-basis-module} above.
Checking the obvious functorialities is left to the reader.
\end{proof}

\noindent
Finally, we address the question of the relationship of this with
continuous maps. This is now very easy thanks to the work
above. First we do the case where there is a basis on the
target given.

\begin{lemma}
\label{lemma-f-map-basis-below-structures}
Let $f : X \to Y$ be a continuous map of topological spaces.
Let $(\mathcal{C}, F)$ be a type of algebraic structures.
Let $\mathcal{F}$ be a sheaf with values in $\mathcal{C}$ on $X$.
Let $\mathcal{G}$ be a sheaf with values in $\mathcal{C}$ on $Y$.
Let $\mathcal{B}$ be a basis for the topology on $Y$.
Suppose given for every $V \in \mathcal{B}$ a morphism
$$
\varphi_V :
\mathcal{G}(V)
\longrightarrow
\mathcal{F}(f^{-1}V)
$$
of $\mathcal{C}$ compatible with restriction mappings.
Then there is a unique $f$-map (see Definition \ref{definition-f-map}
and discussion of $f$-maps in
Section \ref{section-presheaves-structures-functorial})
$\varphi : \mathcal{G} \to \mathcal{F}$
recovering $\varphi_V$ for $V \in \mathcal{B}$.
\end{lemma}

\begin{proof}
This is trivial because the collection of maps
amounts to a morphism between the restrictions
of $\mathcal{G}$ and $f_*\mathcal{F}$ to $\mathcal{B}$.
By Lemma \ref{lemma-restrict-basis-equivalence-structures}
this is the same as giving a morphism from $\mathcal{G}$
to $f_*\mathcal{F}$, which by Lemma \ref{lemma-f-map}
is the same as an $f$-map. See also
Lemma \ref{lemma-f-map-sets-algebraic-structures}
and the discussion preceding it
for how to deal with the case of sheaves of algebraic structures.
\end{proof}

\noindent
Here is the analogue for ringed spaces.

\begin{lemma}
\label{lemma-f-map-basis-below-modules}
Let $(f, f^\sharp) : (X, \mathcal{O}_X) \to (Y, \mathcal{O}_Y)$
be a morphism of ringed spaces.
Let $\mathcal{F}$ be a sheaf of $\mathcal{O}_X$-modules.
Let $\mathcal{G}$ be a sheaf of $\mathcal{O}_Y$-modules.
Let $\mathcal{B}$ be a basis for the topology on $Y$.
Suppose given for every $V \in \mathcal{B}$ a
$\mathcal{O}_Y(V)$-module map
$$
\varphi_V :
\mathcal{G}(V)
\longrightarrow
\mathcal{F}(f^{-1}V)
$$
(where $\mathcal{F}(f^{-1}V)$ has a module structure using
$f^\sharp_V : \mathcal{O}_Y(V) \to \mathcal{O}_X(f^{-1}V)$)
compatible with restriction mappings.
Then there is a unique $f$-map (see discussion of $f$-maps in
Section \ref{section-ringed-spaces-functoriality-modules})
$\varphi : \mathcal{G} \to \mathcal{F}$
recovering $\varphi_V$ for $V \in \mathcal{B}$.
\end{lemma}

\begin{proof}
Same as the proof of the corresponding lemma
for sheaves of algebraic structures above.
\end{proof}

\begin{lemma}
\label{lemma-f-map-basis-above-and-below-structures}
Let $f : X \to Y$ be a continuous map of topological spaces.
Let $(\mathcal{C}, F)$ be a type of algebraic structures.
Let $\mathcal{F}$ be a sheaf with values in $\mathcal{C}$ on $X$.
Let $\mathcal{G}$ be a sheaf with values in $\mathcal{C}$ on $Y$.
Let $\mathcal{B}_Y$ be a basis for the topology on $Y$.
Let $\mathcal{B}_X$ be a basis for the topology on $X$.
Suppose given for every $V \in \mathcal{B}_Y$, and
$U \in \mathcal{B}_X$ such that $f(U) \subset V$ a morphism
$$
\varphi_V^U :
\mathcal{G}(V)
\longrightarrow
\mathcal{F}(U)
$$
of $\mathcal{C}$ compatible with restriction mappings.
Then there is a unique $f$-map (see
Definition \ref{definition-f-map} and the discussion
of $f$-maps in Section \ref{section-presheaves-structures-functorial})
$\varphi : \mathcal{G} \to \mathcal{F}$
recovering $\varphi_V^U$ as the composition
$$
\mathcal{G}(V) \xrightarrow{\varphi_V}
\mathcal{F}(f^{-1}(V)) \xrightarrow{\text{restr.}}
\mathcal{F}(U)
$$
for every pair $(U, V)$ as above.
\end{lemma}

\begin{proof}
Let us first proves this for sheaves of sets.
Fix $V \subset Y$ open. Pick $s \in \mathcal{G}(V)$.
We are going to construct an element
$\varphi_V(s) \in \mathcal{F}(f^{-1}V)$.
We can define a value $\varphi(s)_x$ in the stalk $\mathcal{F}_x$
for every $x \in f^{-1}V$ by picking a $U \in \mathcal{B}_X$
with $x \in U \subset f^{-1}V$ and setting $\varphi(s)_x$
equal to the equivalence class of $(U, \varphi_V^U(s))$
in the stalk. Clearly, the family $(\varphi(s)_x)_{x \in f^{-1}V}$
satisfies condition $(*)$ because the maps $\varphi_V^U$
for varying $U$
are compatible with restrictions in the sheaf $\mathcal{F}$.
Thus, by the proof of Lemma \ref{lemma-extend-off-basis}
we see that $(\varphi(s)_x)_{x \in f^{-1}V}$ corresponds
to a unique element $\varphi_V(s)$ of $\mathcal{F}(f^{-1}V)$.
Thus we have defined a set map
$\varphi_V : \mathcal{G}(V) \to \mathcal{F}(f^{-1}V)$.
The compatibility between $\varphi_V$ and $\varphi_V^U$
follows from Lemma \ref{lemma-condition-star-sections}.

\medskip\noindent
We leave it to the reader to show that the construction
of $\varphi_V$ is compatible with restriction mappings as we vary
$V \in \mathcal{B}_Y$. Thus we may apply Lemma
\ref{lemma-f-map-basis-below-structures} above to
``glue'' them to the desired $f$-map.

\medskip\noindent
Finally, we note that the map of sheaves of sets so constructed
satisfies the property that the map on stalks
$$
\mathcal{G}_{f(x)} \longrightarrow \mathcal{F}_x
$$
is the colimit of the system of maps $\varphi_V^U$ as
$V \in \mathcal{B}_Y$ varies over those elements that
contain $f(x)$ and $U \in \mathcal{B}_X$ varies over those elements that
contain $x$. In particular, if $\mathcal{G}$ and $\mathcal{F}$
are the underlying sheaves of sets of sheaves of algebraic structures,
then we see that the maps on stalks is a morphism of algebraic
structures. Hence we conclude that the associated map of
sheaves of underlying sets $f^{-1}\mathcal{G} \to \mathcal{F}$
satisfies the assumptions of
Lemma \ref{lemma-f-map-sets-algebraic-structures}.
We conclude that $f^{-1}\mathcal{G} \to \mathcal{F}$
is a morphism of sheaves with values in $\mathcal{C}$.
And by adjointness this means that $\varphi$ is
an $f$-map of sheaves of algebraic structures.
\end{proof}

\begin{lemma}
\label{lemma-f-map-basis-above-and-below-modules}
Let $(f, f^\sharp) : (X, \mathcal{O}_X) \to (Y, \mathcal{O}_Y)$
be a morphism of ringed spaces.
Let $\mathcal{F}$ be a sheaf of $\mathcal{O}_X$-modules.
Let $\mathcal{G}$ be a sheaf of $\mathcal{O}_Y$-modules.
Let $\mathcal{B}_Y$ be a basis for the topology on $Y$.
Let $\mathcal{B}_X$ be a basis for the topology on $X$.
Suppose given for every $V \in \mathcal{B}_Y$, and
$U \in \mathcal{B}_X$ such that $f(U) \subset V$ a
$\mathcal{O}_Y(V)$-module map
$$
\varphi_V^U :
\mathcal{G}(V)
\longrightarrow
\mathcal{F}(U)
$$
compatible with restriction mappings. Here the
$\mathcal{O}_Y(V)$-module structure on $\mathcal{F}(U)$
comes from the $\mathcal{O}_X(U)$-module structure
via the map $f^\sharp_V : \mathcal{O}_Y(V)
\to \mathcal{O}_X(f^{-1}V) \to \mathcal{O}_X(U)$.
Then there is a unique $f$-map of sheaves of modules (see
Definition \ref{definition-f-map} and the discussion
of $f$-maps in Section \ref{section-ringed-spaces-functoriality-modules})
$\varphi : \mathcal{G} \to \mathcal{F}$
recovering $\varphi_V^U$ as the composition
$$
\mathcal{G}(V) \xrightarrow{\varphi_V}
\mathcal{F}(f^{-1}(V)) \xrightarrow{\text{restr.}}
\mathcal{F}(U)
$$
for every pair $(U, V)$ as above.
\end{lemma}

\begin{proof}
Similar to the above and omitted.
\end{proof}
